\documentclass[../../main.tex]{subfiles}
\begin{document}

% 年度の大きな表示
\begin{center}
  {\fontsize{48pt}{58pt}\selectfont \textbf{1992}\normalsize\textbf{年度}}
\end{center}

\vspace{10mm}

% 本文


\vspace{15mm}

% 出題テーマと難易度の表
\noindent\textbf{出題テーマと難易度}

\vspace{5mm}

\begin{center}
  \shadowbox{%
    \begin{tabular}{|c|p{8cm}|c|}
      \hline
      \textbf{問題番号} & \textbf{テーマ} & \textbf{難易度} \\
      \hline
      第1問 & 分数関数が整数値をとる条件 & \\
      \hline
      第2問 & 1次変換と不変直線 & \\
      \hline
      第3問 & 放物線と直線で囲まれる面積の最小値 & \\
      \hline
      第4問 & 2進展開と関数の反復合成 & \\
      \hline
      第5問 & 三角関数を含む定積分の漸化式 & \\
      \hline
    \end{tabular}%
  }
\end{center}

\vspace{15mm}

% 解答
\noindent\textbf{解答}

\vspace{5mm}

\begin{center}
  \shadowbox{%
    \renewcommand{\arraystretch}{1.6}
    \begin{tabular}{|c|c|p{8cm}|}
      \hline
      \textbf{問題番号} & \textbf{小問} & \textbf{答え} \\
      \hline
      第1問 & - & $k>3$ \\
      \hline
      第2問 & - & （要確認） \\
      \hline
      第3問 & - & 傾き $m=2$ で最小，最小面積 $\dfrac{4}{3}(c-1)^{3/2}$ \\
      \hline
      \multirow{2}{*}{第4問} & (1) & $f_3(x)=8x-k\ \left(\dfrac{k}{8}\le x<\dfrac{k+1}{8}\right)$（$k=0,1,\cdots,7$） \\
      \cline{2-3}
                             & (2) & $0$ \\
      \hline
      \multirow{2}{*}{第5問} & (1) & $I_n-I_{n-1}=2(-1)^n\left\{\dfrac{\pi}{4(2n-2)}\sin\dfrac{n-1}{2}\pi+\dfrac{1}{(2n-2)^2}\left(\cos\dfrac{n-1}{2}\pi-1\right)\right\}$ \\
      \cline{2-3}
                             & (2) & $I_3=-\dfrac{\pi^2}{32}+\dfrac{\pi}{4}-\dfrac{1}{4}$ \\
      \hline
    \end{tabular}%
    \renewcommand{\arraystretch}{1.0}
  }
\end{center}

\vspace{15mm}

% 解答時間
\noindent\textbf{試験形態}

\vspace{5mm}

\begin{center}
  \shadowbox{%
    \begin{tabular}{p{8cm}}
      （要確認）
    \end{tabular}%
  }
\end{center}

\end{document}
